% GENERATED FILE. Edit canonical lessons, then run npm run generate:books.
% canonical-source-hash: fnv1a64:9ce5ce707d3bcf01
% canonical-lessons: ZH-C86-nar-where, ZH-W86-na-where, ZH-C86-shangbian, ZH-C86-xiabian, ZH-C86-xishoujian, ZH-C86-practice

\chapter{Where Is It?}
\label{ch:zh-where-is-it}
\hlchaptermodality{86}

\begin{chapteropening}
\textbf{By the end of this chapter:} \emph{I can ask where something is and understand the answer.}

Asking where with \zh{哪儿} and placing the answer in the same slot, on top of, under, and the question every traveller needs, built on one new written character: \zh{哪}.
\end{chapteropening}

\section[\texorpdfstring{\zh{哪儿} (nǎr)}{哪儿 (nǎr)}]{\zh{哪儿} (nǎr) — where?}

\label{lesson:ZH-C86-nar-where}



\begin{quote}
A new chapter, and it opens with the question the last one kept answering.

\begin{itemize}
\raggedright
  \item \emph{Recall:} say \zh{我在这儿} — \emph{wǒ zài zhèr}, then \zh{他在那儿} — \emph{tā zài nàr} — \textbf{R1}, the chapter you have finished
  \item \emph{Recall:} say \emph{wàibian}, outside — \textbf{R1}
  \item \emph{Recall:} say \emph{nǐ shì nǎ guó rén}, which country are you from — \textbf{R4}, from the chapter on where people come from
  \item \emph{Recall:} say \emph{liáotiān}, to chat — \textbf{R2}
\end{itemize}
\end{quote}

\subsection*{You'll want to know: \zh{哪儿}}
\textbf{\zh{哪儿}} — \emph{nǎr} — \textquotedblleft{}where?\textquotedblright{}.

\begin{quote}
\textbf{\zh{你在哪儿}} \emph{nǐ zài nǎr} — Where are you?
\end{quote}

\begin{quote}
\textbf{\zh{我在那儿}} \emph{wǒ zài nàr} — I'm over there.
\end{quote}

The next lesson puts the new character in your hand.

\begin{cousinweb}
\textbf{\zh{哪儿}} (\emph{nǎr}) is built from \zh{哪} 'which', the \emph{nǎ} of \emph{nǐ shì nǎ guó rén}, and the r ending \zh{儿}. \zh{哪} has no reading in the Guangyun: no dictionary we use records an older form, so its history is not settled.
\end{cousinweb}

\begin{grammarlens}[title={What you've built}]
The answer goes where the question word was, and nothing else moves: \zh{你在哪儿} — \zh{我在那儿}. Only the tone tells the two words apart: \emph{nǎr}, with the dip, asks; \emph{nàr}, with the fall, answers.
\end{grammarlens}

\subsection*{Your turn}
\begin{itemize}
\raggedright
  \item \emph{Say it:} \emph{nǎr} — the dip
  \item \emph{Say it:} \emph{nàr} — the fall
  \item \emph{Say it:} \zh{你在哪儿} — \emph{nǐ zài nǎr} — Where are you?
  \item \emph{Say it:} \zh{我在那儿} — \emph{wǒ zài nàr} — I'm over there.
  \item \emph{Recall:} say \emph{liáotiān}, then ask where somebody is
\end{itemize}

\begin{tcolorbox}[breakable,colback=teal!4,colframe=teal!35!black,title={Before you move on}]
What does \zh{哪儿} ask? (\textquotedblleft{}Where?\textquotedblright{}.) What tells \emph{nǎr} from \emph{nàr}? (\textbf{The tone: the dip asks, the fall answers}.)
\end{tcolorbox}
\section[\texorpdfstring{\zh{哪} (nǎ)}{哪 (nǎ)}]{\zh{哪} — \zh{口} beside the whole of \zh{那}}

\label{lesson:ZH-W86-na-where}



\begin{quote}
The word first, then the character inside it.

\begin{itemize}
\raggedright
  \item \emph{Recall:} say \emph{nǎr} — \textbf{R1}, one lesson back, before it decays
\end{itemize}

\begin{quote}
\zh{口} \zh{那}
\end{quote}

\emph{Read it:} each one aloud, with nothing said first — \emph{kǒu}, \emph{nà}
\end{quote}

\subsection*{Script — \zh{口} beside the whole of \zh{那}}
\hlblockfigure{\detokenize{figures/ZH-W86-na-where-filmstrip.pdf}}{How \zh{哪} is written, stroke by stroke}

\begin{quote}
\zh{哪}
\end{quote}

\textbf{Two characters you own, side by side.} Nothing in it is new.

\textbf{\zh{口}, three strokes}, small, on the left: the left side, the top that turns down the right side, the bottom bar.

\textbf{\zh{那}, six strokes}, whole, as you write it alone: the hooked turn across the top and down the right, the upper bar, the lower bar, the left-falling stroke through them, then the ear on the right, across, back, round and hooking, and its long vertical.

3 + 6 = nine.

\subsection*{Your turn}
\begin{itemize}
\raggedright
  \item \emph{Write it:} \zh{哪}, all nine, with the model in front of you
  \item \emph{Cover:} write it again from memory, and count the pieces, not the strokes
\end{itemize}

\begin{tcolorbox}[breakable,colback=teal!4,colframe=teal!35!black,title={Before you move on}]
How many strokes in \zh{哪}? (\textbf{Nine}.) Which two characters sit side by side? (\textbf{\zh{口} and \zh{那}}.) What tells \zh{哪} from \zh{那} when you say them? (\textbf{The tone: \emph{nǎ} dips, \emph{nà} falls}.)

Source: \href{https://raw.githubusercontent.com/chanind/hanzi-writer-data/68d10a4b21150cae5e1ebbd223eed289cf32d90c/data/\%E5\%93\%AA.json}{Hanzi Writer Data \zh{哪}, pinned stroke paths and medians}.
\end{tcolorbox}
\section[shàngbian]{shàngbian — on top; above}

\label{lesson:ZH-C86-shangbian}



\begin{quote}
Before the new one: read the question from the last lesson, then two words for things at home.

\begin{itemize}
\raggedright
  \item \emph{Read it:} \zh{你在哪儿}, aloud — \emph{nǐ zài nǎr}
  \item \emph{Recall:} say \emph{zhuōzi}, a table
  \item \emph{Recall:} say \emph{zhōngjiān}, the middle
\end{itemize}
\end{quote}

\subsection*{You'll want to know: shàngbian}
\emph{shàngbian} — \textquotedblleft{}on top; above\textquotedblright{}.

\begin{quote}
\emph{shū zài zhuōzi shàngbian} — The book is on the table.
\end{quote}

Its second character comes later; for now it is a word for the ear.

\begin{cousinweb}
\emph{shàngbian} is built from shàng 'above', the \zh{上} you write, and bian 'side'. Both are inherited from Middle Chinese, \emph{dzyangH} and \emph{pen} (Guangyun, in Baxter's transcription). The Guangyun glosses \emph{dzyangH} as 'ruler', not 'above', so the plain sense 'on top' is assumed here rather than quoted.
\end{cousinweb}

\begin{grammarlens}[title={What you've built}]
A thing, then the place word after it: \emph{zhuōzi shàngbian}, on top of the table. With \zh{在} in front it answers \zh{哪儿}.
\end{grammarlens}

\subsection*{Your turn}
\begin{itemize}
\raggedright
  \item \emph{Say it:} \emph{shàngbian}
  \item \emph{Say it:} \emph{shàngbian}, once more
  \item \emph{Say it:} \emph{shū zài zhuōzi shàngbian} — The book is on the table.
  \item \emph{Say it:} \zh{我的手机在哪儿} — \emph{wǒ de shǒujī zài nǎr} — Where's my phone?
  \item \emph{Recall:} say \emph{zhōngjiān}, then \emph{shàngbian} again
\end{itemize}

\begin{tcolorbox}[breakable,colback=teal!4,colframe=teal!35!black,title={Before you move on}]
What does \emph{shàngbian} mean? (\textquotedblleft{}On top; above\textquotedblright{}.) Which written character does it open with? (\textbf{\zh{上}}.)
\end{tcolorbox}
\section[xiàbian]{xiàbian — under; below}

\label{lesson:ZH-C86-xiabian}



\begin{quote}
Before the new one: last lesson's word, then two words you have had for a while.

\begin{itemize}
\raggedright
  \item \emph{Recall:} say \emph{shàngbian}, on top — \textbf{R1}, one lesson back
  \item \emph{Recall:} say \emph{yǐzi}, a chair — \textbf{R3}, durable
  \item \emph{Recall:} say \emph{duìmiàn}, opposite — \textbf{R2}
\end{itemize}
\end{quote}

\subsection*{You'll want to know: xiàbian}
\emph{xiàbian} — \textquotedblleft{}under; below\textquotedblright{}.

\begin{quote}
\emph{māo zài yǐzi xiàbian} — The cat is under the chair.
\end{quote}

Its second character comes later; for now it is a word for the ear.

\begin{cousinweb}
\emph{xiàbian} is built from xià 'below', the \zh{下} you write, and bian 'side'. Both are inherited from Middle Chinese, \emph{haeX} and \emph{pen} (Guangyun, in Baxter's transcription), \emph{haeX} glossed 'to go down' and 'the bottom'.
\end{cousinweb}

\begin{grammarlens}[title={What you've built}]
\emph{shàngbian} and \emph{xiàbian}, on top and under: the pair you write as \zh{上} and \zh{下}, with the same ending as \emph{lǐbian} and \emph{wàibian}.
\end{grammarlens}

\subsection*{Your turn}
\begin{itemize}
\raggedright
  \item \emph{Say it:} \emph{xiàbian}
  \item \emph{Say it:} \emph{xiàbian}, once more
  \item \emph{Say it:} \emph{māo zài yǐzi xiàbian} — The cat is under the chair.
  \item \emph{Say it:} \emph{shū zài zhuōzi shàngbian} — The book is on the table.
  \item \emph{Recall:} say \emph{duìmiàn}, then \emph{xiàbian} again
\end{itemize}

\begin{tcolorbox}[breakable,colback=teal!4,colframe=teal!35!black,title={Before you move on}]
What does \emph{xiàbian} mean? (\textquotedblleft{}Under; below\textquotedblright{}.) And on top? (\emph{\textbf{shàngbian}}.)
\end{tcolorbox}
\section[xǐshǒujiān]{xǐshǒujiān — toilet; washroom}

\label{lesson:ZH-C86-xishoujian}



\begin{quote}
Before the new one: last lesson's word, then the polite opener you will need with it.

\begin{itemize}
\raggedright
  \item \emph{Recall:} say \emph{xiàbian}, under — \textbf{R1}, one lesson back
  \item \emph{Recall:} say \zh{请问} — \emph{qǐngwèn}, excuse me, may I ask
  \item \emph{Recall:} say \emph{fángjiān}, a room
\end{itemize}
\end{quote}

\subsection*{You'll want to know: xǐshǒujiān}
\emph{xǐshǒujiān} — \textquotedblleft{}toilet; washroom\textquotedblright{}.

\begin{quote}
\emph{qǐngwèn, xǐshǒujiān zài nǎr} — Excuse me, where is the toilet?
\end{quote}

Its first and last characters come later; for now it is a word for the ear.

\begin{cousinweb}
\emph{xǐshǒujiān} is built from xǐ 'to wash', shǒu 'hand', the \zh{手} you write, and jiān 'room', the jiān of \emph{fángjiān}. All three are inherited from Middle Chinese, \emph{sejX}, \emph{syuwX} and \emph{kean} (Guangyun, in Baxter's transcription). The Guangyun glosses \emph{sejX} 'to wash' and \emph{kean} 'a gap; between'; 'a room' is a later sense of jiān.
\end{cousinweb}

\begin{grammarlens}[title={What you've built}]
\emph{qǐngwèn, xǐshǒujiān zài nǎr} may be the most useful sentence in the chapter: a polite opener, the thing, \zh{在}, and the question word in the place slot.
\end{grammarlens}

\subsection*{Your turn}
\begin{itemize}
\raggedright
  \item \emph{Say it:} \emph{xǐshǒujiān}
  \item \emph{Say it:} \emph{xǐshǒujiān}, once more
  \item \emph{Say it:} \emph{qǐngwèn, xǐshǒujiān zài nǎr} — Excuse me, where is the toilet?
  \item \emph{Recall:} say \emph{xiàbian} and \emph{fángjiān}, then \emph{xǐshǒujiān} again
\end{itemize}

\begin{tcolorbox}[breakable,colback=teal!4,colframe=teal!35!black,title={Before you move on}]
What does \emph{xǐshǒujiān} mean? (\textquotedblleft{}Toilet; washroom\textquotedblright{}.) Ask where it is, politely. (\emph{qǐngwèn, xǐshǒujiān zài nǎr}.)
\end{tcolorbox}
\section[Practice]{\zh{你在哪儿} — \zh{我在那儿}}

\label{lesson:ZH-C86-practice}



\begin{quote}
The chapter closes here, with a phone that is not where it should be.

\begin{itemize}
\raggedright
  \item \emph{Recall:} say \emph{xǐshǒujiān}, the toilet — \textbf{R1}, one lesson back
  \item \emph{Recall:} say \emph{cānjiā}, to take part — \textbf{R2}
  \item \emph{Recall:} say \emph{lǐbian}, inside — \textbf{R2}, from the chapter before
\end{itemize}
\end{quote}

\subsection*{The exchange}
\begin{quote}
\textbf{A:} \zh{我的手机呢}
\end{quote}

\begin{quote}
\emph{wǒ de shǒujī ne} — Where's my phone?
\end{quote}

\begin{quote}
\textbf{B:} \emph{zài zhuōzi shàngbian} — On the table.
\end{quote}

\begin{quote}
\textbf{A:} \emph{bú zài zhuōzi shàngbian} — It isn't on the table.
\end{quote}

\begin{quote}
\textbf{B:} \emph{zài yǐzi xiàbian} — Under the chair, then.
\end{quote}

\begin{quote}
\textbf{A:} \zh{对} \zh{谢谢}
\end{quote}

\begin{quote}
\emph{duì, xièxie} — Yes. Thanks.
\end{quote}

\begin{quote}
\textbf{A:} \emph{qǐngwèn, xǐshǒujiān zài nǎr} — Excuse me, where's the toilet?
\end{quote}

\begin{quote}
\textbf{B:} \zh{在那儿}
\end{quote}

\begin{quote}
\emph{zài nàr} — Over there.
\end{quote}

\begin{quote}
\textbf{A:} \zh{你在哪儿}
\end{quote}

\begin{quote}
\emph{nǐ zài nǎr} — And where will you be?
\end{quote}

\begin{quote}
\textbf{B:} \emph{wǒ zài wàibian} — Outside.
\end{quote}

Every answer starts at \zh{在}, or drops even that. The thing being looked for was named once, and the rest of the exchange leaves it out.

\begin{tcolorbox}[breakable,skin=enhanced,colback=orange!6,colframe=orange!45!black,boxrule=0.5pt,arc=1mm,left=6pt,right=6pt,top=4pt,bottom=4pt,fonttitle=\bfseries,title={Read this}]
\emph{Read it:} the exchange through once without stopping, then answer the question below
\end{tcolorbox}

\begin{grammarlens}[title={What you've built — one question, any answer}]
One question, \textbf{\zh{你在哪儿}} \emph{nǐ zài nǎr}, and the answer slots into the same place:

\begin{itemize}
\raggedright
  \item \textbf{\zh{我在那儿}} \emph{wǒ zài nàr} — over there.
  \item \emph{zài zhuōzi shàngbian} — on the table.
  \item \emph{zài yǐzi xiàbian} — under the chair.
\end{itemize}

\textbf{Listen.} \emph{Nǐ de shǒujī zài nǎr?}

\textbf{Speak.} Answer it three ways, then ask it back.

\emph{Write it:} \zh{哪} — \zh{口} first, then \zh{那}
\end{grammarlens}

\subsection*{Your turn}
\begin{itemize}
\raggedright
  \item \emph{Say it:} \emph{qǐngwèn, xǐshǒujiān zài nǎr} — Excuse me, where's the toilet?
  \item \emph{Say it:} \zh{我的手机呢} — \emph{wǒ de shǒujī ne} — Where's my phone?
  \item \emph{Say it:} the whole exchange, both parts
  \item \emph{Write it:} \zh{你在哪儿}
\end{itemize}

\begin{tcolorbox}[breakable,colback=teal!4,colframe=teal!35!black,title={Before you move on}]
Ask somebody where they are. (\textbf{\zh{你在哪儿}}.) Ask politely where the toilet is. (\emph{qǐngwèn, xǐshǒujiān zài nǎr}.) Which tone asks, \emph{nǎr} or \emph{nàr}? (\textbf{The dip, \emph{nǎr}}.)
\end{tcolorbox}
